% ભાગ: first-order-logic
% પ્રકરણ: beyond
% વિભાગ: many-sorted-logic

\documentclass[../../../include/open-logic-section]{subfiles}

\begin{document}

\olfileid{fol}{byd}{msl}

\olsection{બહુ-પ્રકાર તર્કશાસ્ત્ર}

પ્રથમ-ક્રમ તર્કશાસ્ત્રમાં ચલો અને પરિમાણકો એક જ !!{domain} પર
વ્યાપે છે. પરંતુ ઘણી વાર અનેક (પરસ્પર વિયુક્ત) !!{domain}s રાખવા
ઉપયોગી બને છે: દાખલા તરીકે, સંખ્યાઓનું !!a{domain}, ભૌમિતિક વસ્તુઓનું
!!a{domain}, સંખ્યાઓથી સંખ્યાઓમાં જતાં વિધેયોનું !!a{domain},
ક્રમવિનિમેય સમૂહોનું !!a{domain}, વગેરે રાખવા આપણે માગી શકીએ.

બહુ-પ્રકાર તર્કશાસ્ત્ર આવું માળખું આપે છે. શરૂઆત ``પ્રકારો''ની એક
યાદીથી થાય છે---કોઈ વસ્તુનો ``પ્રકાર'' તે કયા ``!!{domain}''માં રહેવાની
છે તે દર્શાવે છે. પછી દરેક પ્રકાર માટે !!{variable}s અને પરિમાણકો, અને
(સામાન્ય રીતે) દરેક પ્રકાર માટે !!{identity} હોય છે. વિધેયો અને સંબંધો
પણ તેઓ દલીલો તરીકે લઈ શકે એવી વસ્તુઓના પ્રકારો વડે ``પ્રકારિત'' હોય
છે. બાકી પ્રથમ-ક્રમ તર્કશાસ્ત્રના સામાન્ય નિયમો જ રાખવામાં આવે છે અને
દરેક પ્રકાર માટે પરિમાણક-નિયમોનાં સ્વરૂપો પુનરાવર્તિત થાય છે.

દાખલા તરીકે, આંતરરાષ્ટ્રીય સંબંધોના અભ્યાસ માટે આપણે વસ્તુઓના બે
પ્રકારો---ફ્રેન્ચ નાગરિકો અને જર્મન નાગરિકો---ધરાવતી ભાષા પસંદ કરી
શકીએ. પહેલા પ્રકારની વસ્તુઓ માટે ``વાઇન પીવે છે'' એવો એકસ્થાની સંબંધ;
બીજા પ્રકારની વસ્તુઓ માટે ``વુર્સ્ટ ખાય છે'' એવો બીજો એકસ્થાની સંબંધ;
અને બે દલીલો લેતો ``બહુરાષ્ટ્રીય વિવાહિત યુગલ બનાવે છે'' એવો દ્વિસ્થાની
સંબંધ રાખી શકીએ, જેમાં પહેલી દલીલ પહેલા પ્રકારની અને બીજી દલીલ બીજા
પ્રકારની હોય. ફ્રેન્ચ નાગરિકો પર વ્યાપવા માટે $a$, $b$, $c$ અને જર્મન
નાગરિકો પર વ્યાપવા માટે $x$, $y$, $z$ ચલો વાપરીએ, તો
\[
\lforall[a][\lforall[x][(\Atom{\Obj{MarriedTo}}{a,x} \lif
(\Atom{\Obj{DrinksWine}}{a} \lor \lnot \Atom{\Obj{EatsWurst}}{x})]]
\]
કહે છે કે કોઈ ફ્રેન્ચ વ્યક્તિ કોઈ જર્મન વ્યક્તિ સાથે વિવાહિત હોય, તો
કાં તો ફ્રેન્ચ વ્યક્તિ વાઇન પીવે છે અથવા જર્મન વ્યક્તિ વુર્સ્ટ ખાતી
નથી.
\sourcecorrection{OLBYD-001}{મૂળની
\texttt{many-sorted-logic.tex}, પંક્તિ~38માં બીજા સાર્વત્રિક પરિમાણકને
\texttt{\string\lforall x]} તરીકે લખવામાં આવ્યો છે અને તેનો આરંભક
ચોરસ કૌંસ ખૂટે છે. અહીં સુમેળવાળું~\texttt{\string\lforall[x]}
પુનઃસ્થાપિત કર્યું છે.}

બહુ-પ્રકાર તર્કશાસ્ત્રને કુદરતી રીતે પ્રથમ-ક્રમ તર્કશાસ્ત્રમાં સમાવી
શકાય છે: બહુ-પ્રકાર !!{domain}sની બધી વસ્તુઓને એક પ્રથમ-ક્રમ
!!{domain}માં એકત્ર કરી, પ્રકારોનો હિસાબ રાખવા એકસ્થાની !!{predicate}s
વાપરી અને પરિમાણકોને સાપેક્ષ બનાવવામાં આવે છે. દાખલા તરીકે, ઉપરના
ઉદાહરણને અનુરૂપ પ્રથમ-ક્રમ ભાષામાં વર્ણવેલા બીજા સંબંધો ઉપરાંત
``$\Obj{German}$'' અને ``$\Obj{French}$'' એવા એકસ્થાની !!{predicate}s
હશે, અને પ્રકારની શરતો દૂર કરેલી હશે. જર્મન પ્રકારનો !!a{variable}
$x$ ધરાવતો પ્રકારિત પરિમાણક $\lforall[x][!A]$ આ રીતે અનુવાદિત થાય છે:
\[
\lforall[x][(\Atom{\Obj{German}}{x} \lif !A)].
\]
પ્રકારો અલગ છે તેની ખાતરી કરાવતી સ્વયંસિદ્ધિઓ---દાખલા તરીકે,
$\lforall[x][\lnot (\Atom{\Obj{German}}{x} \land
  \Atom{\Obj{French}}{x})]$---તેમજ ``વાઇન પીવે છે'' માત્ર
$\Atom{\Obj{French}}{x}$ વિધેય સંતોષતી વસ્તુઓ માટે જ લાગુ પડે છે તેની
ખાતરી કરાવતી સ્વયંસિદ્ધિઓ વગેરે ઉમેરવી પડે. આ રૂઢિઓ અને સ્વયંસિદ્ધિઓ
સાથે બહુ-પ્રકાર !!{sentence}s પ્રથમ-ક્રમ !!{sentence}sમાં અને
બહુ-પ્રકાર !!{derivation}s પ્રથમ-ક્રમ !!{derivation}sમાં અનુવાદિત થાય
છે એવું બતાવવું મુશ્કેલ નથી. વળી, બહુ-પ્રકાર !!{structure}s અનુરૂપ
પ્રથમ-ક્રમ !!{structure}sમાં અને ઊલટું ``અનુવાદિત'' થાય છે; તેથી
બહુ-પ્રકાર તર્કશાસ્ત્ર માટે પણ આપણી પાસે પૂર્ણતા પ્રમેય છે.

\end{document}
